\selectlanguage{ngerman}%
\frenchspacing%
\begin{ledgroupsized}[r]{120mm}%
\footnotesize%
\pstart%
\noindent%
\textbf{Überlieferung:}%
\pend%
\end{ledgroupsized}%
\begin{ledgroupsized}[r]{114mm}%
\footnotesize%
\pstart%
\parindent -6mm%
\makebox[6mm][l]{\textit{L}}%
Aufzeichnung:
LH XXXVII 4 Bl.\;104.
Ein Blatt 4\textsuperscript{o},
Ränder beschnitten:
Papiererhaltungsmaßnahmen;
Fragment eines
\href{\Scan{LH_37_04_104wz}}{Wasserzeichens}:
Harzer Papier. %
Eineinhalb Seiten;
untere Hälfte von Bl.\;\href{\Scan{LH_37_04_104v}}{104\,v\textsuperscript{o}} unbeschrieben.
\pend%
\end{ledgroupsized}%
\vspace{5mm}%
\begin{ledgroup}%
\footnotesize%
\pstart%
\noindent%
\textbf{Datierungsgründe:}
Die Datierung der Aufzeichnung N.\;\ref{60258} beruht auf äußeren Eigenschaften des Texträgers.
Siehe die ausführliche Begründung auf S.\;\pageref{60258_dat}.%
\pend%
\end{ledgroup}%
\selectlanguage{latin}%
\frenchspacing%
\vspace*{8mm}
\normalsize%
\pstart%
\noindent%
\lbrack%
\href{\Scan{LH_37_04_104r}}%
{104\,r\textsuperscript{o}}%
\rbrack%
\hfill%
Observationes von Winden%
\indexS{observatio}%
\indexS{Wind}%
\hfill%
\phantom{\lbrack104\,r\textsuperscript{o}\rbrack}%
\pend%
\pstart%
\noindent%
Wenn ein sud Wind%
\indexS{Südwind}
oder Sud westwind%
\indexS{Südwestwind}
etliche tage geblasen,%
\indexS{Tag}
und ein Nord%
\indexS{Nordwind}
oder Nordost wind%
\indexS{Nordostwind}
folget,
erhebt
\edtext{%
sich der wetter Zeiger%
\indexS{Wetterzeiger}%
\indexS{Zeiger}%
}{%
\lemma{sich}%
\Bfootnote{%
\textit{(1)}~der sudwest wind im wet %
\textit{(2)}~der wetter Zeiger%
~\textit{L}%
}}
7 oder 8 linien höher
als er geweßen,
und steigt etwa auff 28 Zoll,%
\indexS{Zoll}
oder 28 Zoll und etliche linien,
und da folgt sehr schohn wetter.%
\indexS{Wetter}%
\pend%
\pstart%
Wenn nach
\edtext{%
einem Nord%
\indexS{Nordwind}
oder Nordostwind%
\indexS{Nordostwind}
ein Sud%
\indexS{Südwind}
oder Sud West Wind%
\indexS{Südwestwind}%
}{%
\lemma{einem}%
\Bfootnote{%
\textit{(1)}~sud oder sud west Winde ein %
\textit{(2)}~Nord oder %
\lbrack...\rbrack\ West Wind%
~\textit{L}%
}}
folget,
fallet der wetter Zeiger%
\indexS{Wetterzeiger}%
\indexS{Zeiger}
etwa auff 27 Zoll%
\indexS{Zoll}
\edtext{%
und 4 linien%
}{%
\lemma{und 4}%
\Bfootnote{%
\textit{(1)}~ligne %
\textit{(2)}~linien%
~\textit{L}%
}}
oder
\edtext{%
bis weilen%
}{%
\lemma{bis weilen}%
\Bfootnote{%
\textit{erg.\;L}%
}}
27 Zoll,
oder gar 26 Zoll%
\indexS{Zoll}
10 linien.
Und da folget großer regen.%
\indexS{Regen}
Es beträgt sich bis weilen
daß sud und Sudwest viel wolcken%
\indexS{Wolken}
nacher Nord und Nordost getrieben,
und eine Ebbe%
\indexS{Ebbe}
oder ruckfluß%
\indexS{Rückfluss der Luft}
der lufft%
\indexS{Lufft, zurückfließend}
erfolget,
da macht der Nord%
\indexS{Nordwind}
oder NO,%
\indexS{Nordostwind}
daß die zuruck getriebenen Wolcken%
\indexS{Wolken}
zusammen gepreßet werden
und ein oder 2 tage%
\indexS{Tag}
regen folget.%
\indexS{Regen}%
\pend%
\pstart%
Wenn Nord Winde%
\indexS{Nordwind}
aufhohren%
\lbrack,\rbrack\
pflegen ost winde%
\indexS{Ostwind}
zu folgen,
und denn sud%
\indexS{Südwind}
und sud West,%
\indexS{Südwestwind}
und selten contra.%
\pend%
\pstart%
Deßen ursachen%
\indexS{Ursache}
ich mir also einbilde:
weil die erde%
\indexS{Erde}
umb ihren ax herumb
\edtext{%
sich drehet%
}{%
\lemma{sich}%
\Bfootnote{%
\textit{(1)}~gedrehet %
\textit{(2)}~drehet%
~\textit{L}%
}}%
\lbrack,\rbrack\
folgt ihr die lufft%
\indexS{Lufft, der Erdbewegung folgend}
so nahe bey der lini,
doch etwas langsamer%
\lbrack;\rbrack\
dahehr alda meist ostwind%
\indexS{Ostwind}
gegen die Tropicos%
\indexS{Tropicus}
und 10 oder 12 grad drüber,
da ist die lufft%
\indexS{Lufft, langsamer}
viel langsamer als unter der lini%
\lbrack;\rbrack\
dahehr zu zeiten ein theil der erwarmten lufft%
\indexS{Lufft, erwärmt}
von der lini nach den Tropicos gehet,%
\indexS{Tropicus}
daher ein Sud West%
\indexS{Südwestwind}
oder West%
\indexS{Westwind}%
\lbrack:\rbrack\
gleichwie das waßer%
\indexS{Waßer}%
\lbrack,\rbrack\
so mitten in strohm schnell,%
\indexS{Schtromschnell}
wellen wirfft,%
\indexS{Welle}
welche etwas scharff nach dem ufer%
\indexS{Ufer}
zu gehen,
bis weilen auch wohl rückwerts.
Wenn dieße wellen nachlaßen,%
\indexS{Welle}
und die Nord%
\indexS{Nordwind}
und nordost winde%
\indexS{Nordostwind}
ihre Ebbe%
\indexS{Ebbe}
verrichtet haben%
\lbrack,\rbrack\
pflegets stille zu werden,
und folgt ein ostwind%
\indexS{Ostwind}
so von Natur aus%
\indexS{Natur}
\edtext{%
\lbrack der\rbrack%
}{%
\lemma{der}%
\Bfootnote{%
\textit{erg.\;Hrsg.}%
}}
bewegung der Erde%
\indexS{Bewegung der Erde}%
\indexS{Erde}
ohne dem folget.
Sud%
\indexS{Südwind}
und sudwest%
\indexS{Südwestwind}
folgen auf ost%
\indexS{Ostwind}
in den temperirten Zonis,%
\indexS{Zone, temperiert}
denn es kommen wieder neue wellen%
\indexS{Welle}
von der lini.
Doch wird die ordnung%
\indexS{Ordnung}
bisweilen von exhalation,%
\indexS{Exhalation}
mond,%
\indexS{Mond}
warme%
\indexS{Wärme}
und kalte%
\indexS{Kälte}
und dergleichen particularibus geändert.%
\indexS{particularia}
Nord%
\indexS{Nordwind}
und NO%
\indexS{Nordostwind}
heben den Zeiger%
\indexS{Zeiger}%
\indexS{Wetterzeiger}
denn sie machen
\edtext{%
nicht nur%
}{%
\lemma{nicht nur}%
\Bfootnote{%
\textit{erg.\;L}%
}}
die lufft schwehrer%
\indexS{Lufft, schwerer}
weil sie sie zusamen preßen%
\indexS{Lufft, zusammengepresst}
(\Denarius)
sondern in dem sie sie preßen gegen die Erde,%
\indexS{Erde}
vermehren sie ihr Ressort.%
\indexS{Ressort der Luft}
Und weil NO%
\indexS{Nordostwind}
guth wetter bringet,%
\indexS{Wetter}
so ist die Elevation%
\indexS{Elevation des Zeigers}
des Zeigers%
\indexS{Wetterzeiger}%
\indexS{Zeiger}
ein Zeichen%
\indexS{Zeichen}%
\indexS{Wetterzeichen}
daß Nord%
\indexS{Nordwind}
von oben nieder blase%
\edtext{%
\lbrack; das\rbrack%
}{%
\lemma{; das}%
\Bfootnote{%
\textit{erg.\; Hrsg.}%
}}
bestarcket diese erfahrung,%
\indexS{Erfahrung}
henge eine bleyerne Kugel%
\indexS{Kugel, bleiern}
von~3 Zoll%
\indexS{Zoll}
an einen faden,%
\indexS{Faden}
und gieb ihr eine bewegung%
\indexS{Bewegung in die Runde}
in die rundte%
\indexS{Runde}
in einem eymer%
\indexS{Eimer}%
\indexS{Waßereimer}
waßer,%
\indexS{Waßer}
so wird der staub%
\indexS{Staub}%
\lbrack,\rbrack\
so am boden des eymers%
\indexS{Eimer}%
\indexS{Waßereimer}%
\lbrack,\rbrack\
sich erheben gegen die kugel%
\lbrack;\rbrack\
wenn sie aber~3 oder~4 Zoll%
\indexS{Zoll}%
\lbrack,\rbrack\
nicht davon.
Aber das waßer%
\indexS{Waßer}%
\lbrack,\rbrack\
so nächst bey der kugel%
\indexS{Kugel, bleiern}%
\lbrack,\rbrack\
drehet sich mit ihr umb.
\pend%
\pstart%
\edlabel{LH_37_04_104_paradox-1}%
N und NO%
\indexS{Nordwind}%
\indexS{Nordostwind}
bringen guth wetter erst,%
\indexS{Wetter}
weil sie uber lauter land%
\indexS{Land}
kommen%
\lbrack;\rbrack\
denn
weil sie niederwerts blaßen,
und die elevation%
\indexS{Elevation der Dünste}
der dunste%
\indexS{Dunst}
aus der erde%
\indexS{Erde}
verhindern%
\lbrack;\rbrack\
drittens weil in gepreßeter%
\indexS{Lufft, gepreßt}
\edtext{%
lufft %
der dunst%
\indexS{Dunst}%
}{%
\lemma{lufft}%
\Bfootnote{%
\textit{(1)}~das waßer %
\textit{(2)}~der dunst%
~\textit{L}%
}}
\lbrack%
\href{\Scan{LH_37_04_104v}}%
{104\,v\textsuperscript{o}}%
\rbrack\
beßer schwebet,
\edtext{%
und nicht so%
}{%
\lemma{und}%
\Bfootnote{%
 \textit{(1)}~so %
\textit{(2)}~nicht so%
~\textit{L}%
}}
leicht fallet
und zu tropfen wird.%
\indexS{Tropfen}%
\edlabel{LH_37_04_104_paradox-2}
Der ost wind%
\indexS{Ostwind}
macht nebel%
\indexS{Nebel}
sonderlich im winter%
\indexS{Winter}
und die andern winde%
\indexS{Wind}
\edtext{%
selten,
weil seine bewegung%
\indexS{Bewegung des Windes}%
}{%
\lemma{selten,}%
\Bfootnote{%
\textit{(1)}~weil er %
\textit{(2)}~weil seine bewegung%
~\textit{L}%
}}
nicht auf- noch niederwärts gehet,
sondern der erde parallel,%
\indexS{Erde}
gegen eine langsamere lufft,%
\indexS{Lufft, langsamer}
also bleiben auch die vapores%
\indexS{vapor}
bey ihrer ehemahligen höhe,
und werden von andern vermehret,
und dicker.%
\pend%
\pstart%
Sud und sudwest%
\indexS{Südwind}%
\indexS{Südwestwind}
so von ferne kommen,
blaßen nach den tangent-linien%
\indexS{Tangentlinie}
der Erde%
\indexS{Erde}
und erheben also die hohere lufft,%
\indexS{Lufft, höher}
dadurch wird die gespante feder%
\indexS{Feder, gespannt}%
\indexS{Feder der Lufft}
der untren
lufft%
\indexS{Lufft, untere}
losgelaßen, %
und der Zeiger fallet,%
\indexS{Zeiger}%
\indexS{Wetterzeiger}
da folget regen,%
\indexS{Regen}
sonderlich %
wenn der wind%
\indexS{Wind}
von west %
nach sud und sudwest zuruckgehet,
aber wenn er
\edtext{%
von ONO%
}{%
\lemma{von}%
\Bfootnote{%
\textit{(1)}~N %
\textit{(2)}~ONO%
~\textit{L}%
}}
nach Nord oder NNO gehet,
ist ein Zeichen%
\indexS{Zeichen}%
\indexS{Wetterzeichen}
\edtext{%
da\lbrack ß\rbrack%
}{%
\lemma{das}%
\Bfootnote{%
\textit{L\;ändert Hrsg.}%
}}
guth wetter%
\indexS{Wetter}
folgen wird,
wenn gleich der zeiger%
\indexS{Zeiger}%
\indexS{Wetterzeiger}
ein wenig fallen solte.
Die winde%
\indexS{Wind}
in diesen Landen%
\indexS{Land}
gehen gemeiniglich
von O nach S oder SW
denn W, N, NO,
und
\edtext{%
gar%
}{%
\lemma{gar}%
\Bfootnote{%
\textit{erg.\;L}%
}}
selten gehen sie also contra,
daß sie
\edtext{%
de\lbrack n\rbrack%
}{%
\lemma{denn}%
\Bfootnote{%
\textit{L\;ändert Hrsg.}%
}}
ganzen umbkreiß%
\indexS{Umkreis}
machen solten.%
\pend%
\pstart%
\edlabel{LH_37_04_104_paradox-3}%
Das fallen des zeigers%
\indexS{Zeiger}%
\indexS{Wetterzeiger}
bedeutet regen,%
\indexS{Regen}
erst weil denn die lufft dünner%
\indexS{Lufft, dünner}
und alsdann die dünste%
\indexS{Dunst}
nicht so wohl träget,
da fallen die obersten auf die untersten
und machen große wolcken,%
\indexS{Wolken}
daraus endtlich regen wird.%
\indexS{Regen}%
\edlabel{LH_37_04_104_paradox-4}
Vors andere so kommen sued und sudwest%
\indexS{Südwind}%
\indexS{Südwestwind}
über meer,%
\indexS{Meer}
und sind also mit feuchtigkeiten beladen.%
\indexS{Feuchtigkeit}%
\pend%
\pstart%
Wenn N und NO lang regiret,%
\indexS{Nordwind}%
\indexS{Nordostwind}
fället der Zeiger allmahlig,%
\indexS{Zeiger}%
\indexS{Wetterzeiger}
und doch bleibt guth wetter,%
\indexS{Wetter}
denn die winde%
\indexS{Wind}
bringen wenig dünste%
\indexS{Dunst}
aber die lufft%
\indexS{Lufft, gepreßt}%
\lbrack,\rbrack\
so nieder gepreßet, %
breitet sich auff die seite aus,
und erlanget ihre freyheit wieder.%
\indexS{Freiheit der Luft}
\pend%
\vspace{3.0em}%
\pstart%
\noindent%
\centering%
*     *     *
\pend%
\vspace{2.0em}%
\selectlanguage{german}%
\frenchspacing%
\footnotesize%
\pstart%
\noindent%
\textls{Ausführliche Datierungsgründe:}%
\pagelabel{60258_dat}
Im Textträger der vorliegenden undatierten Aufzeichnung N.\;\ref{60258} \textit{Observationes von Winden} liegt das \href{\Scan{LH_37_04_104wz}}{Fragment} der Gegenmarke eines häufigen Wasserzeichens, das zur Papiermühle bei Arensburg\indexO{Arensburg (Niedersachsen)} gehörte. %
Diese Papiersorte ist im Nachlass nach heutigem Kenntnisstand für die Jahre 1683--1684 belegt.
Mangels weiterer Anhaltspunkte wird zur chronologischen Einordnung von N.\;\ref{60258} die Datierung des Wasserzeichens übernommen.
Die Aufzeichnung könnte aber auch etwas früher oder etwas später entstanden sein.
\pend%
\pstart%
Es gilt zu bemerken, dass in der Aufzeichnung N.\;\ref{60258} Leibniz auch die Frage streift, warum der Luftdruck bei schönem Wetter steigt und bei schlechtem Wetter sinkt;
vgl. \refpassage{LH_37_04_104_paradox-1}{LH_37_04_104_paradox-2} und vor allem \refpassage{LH_37_04_104_paradox-3}{LH_37_04_104_paradox-4}.
Die an diesen Stellen angedeutete Antwort unterscheidet sich von der Erklärung, die er mehr als zehn Jahre später im Textkomplex N.\;\ref{paradoxon_intro} \textit{Paradoxon de gravitate aeris} entfalten würde.
\pend%
\selectlanguage{latin}%
\frenchspacing%
\normalsize%
\newpage%
